\subsection{Spinal Inverse Dynamics With External Loads}
\label{sec:spn-10a}

This section documents \texttt{inverse\_dynamics.py} (issue \#12760, epic
\#12735). It computes the generalized effort of every spinal level, and the
pelvis residual wrench, from a prescribed motion with OpenSim
\texttt{InverseDynamicsSolver}.

\paragraph{Observable.}
The output is \texttt{LoadObservable.GENERALIZED\_EFFORT} (ADR-0055 item 5):
the moment exerted on the superior body by the inferior body through the three
rotational coordinates of a level, expressed in the inferior vertebra ISB
frame ($X$ anterior, $Y$ cranial, $Z$ right) about the mid-disc joint centre,
in N$\cdot$m. It is not a joint reaction. The force array is exactly zero:
the level joints have constant zero translations, so no joint force is an
unknown of the inverse-dynamics problem. Passive element loads and joint
reactions are separate observables and are not mixed in.

\paragraph{Couplers.}
Couplers prescribe kinematics only and are disabled for inverse dynamics
(ADR-0055 item 6): every constraint of the loaded model copy has
\texttt{isEnforced} set to false, which frees every spinal coordinate. The
provenance key \texttt{couplers} records the number disabled. A massless
coupling driver body, if present, contributes zero generalized force.

\paragraph{Muscles.} Every muscle is switched off in the run model, so the
generalized effort is the net joint moment that the muscles, passive elements
and joint reactions must jointly supply; the provenance key \texttt{muscles}
records the exclusion. Muscle forces enter only through static optimization
(SPN-10b) and joint reactions (SPN-10c).

\paragraph{Generalized Force to Moment.}
With ISB angles $\theta_k = s_k q_k$ ($s_k=\pm1$ the inventory axis sign) in
ZXY order, the angular-velocity columns in the parent offset frame $P$ of the
joint are
\[
  J(\theta) = \bigl[\,e_z,\; R_z(\theta_1)e_x,\; R_z(\theta_1)R_x(\theta_2)e_y\,\bigr],
  \qquad \det J = \cos\theta_2 .
\]
Power balance $\tau_\theta = J^{\mathsf T} M_P$ with $\tau_{\theta,k}=s_k\tau_k$
gives $M_P = J^{-\mathsf T}\tau_\theta$, and the moment in the inferior ISB
frame is $M = R_{i0}\,M_P$ with $R_{i0}$ the neutral-offset rotation of the
inventory. A level whose lateral bending is within $10^{-6}$ of $\pm 90^\circ$
raises \texttt{ValueError}. Because the solver returns the generalized force
that must act to produce the motion, a static pose yields the moment that
supports gravity on all bodies above the level.

\paragraph{Pelvis Residual.}
The six generalized forces of the ground-pelvis joint are the residual wrench
about the pelvis origin in the ground frame. For a free joint (quaternion
state) OpenSim speeds are the ground-frame angular velocity and the origin
velocity, so the generalized force is the wrench; a runtime probe verifies
this convention for the model. For a custom joint with Euler coordinates
$M = J^{-\mathsf T}\tau_{\mathrm{rot}}$ with $J$ of the pelvis sequence. Peak
and RMS of the force and moment norms are reported, with
\[
  r_F = \frac{\max\|F\|}{m g},\qquad
  r_M = \frac{\max\|M\|}{m g\,h},
\]
judged against a declared \texttt{ResidualBudget} (defaults $r_F\le 0.05$ and
$r_M\le 0.01$). The height $h$ is the option \texttt{body\_height\_m} or, if
absent, the extent of the body origins along the model's own gravity direction
$\hat g = g/\|g\|$ at the \emph{first} frame of the motion,
$h=\max_b \hat g\!\cdot\!p_b - \min_b \hat g\!\cdot\!p_b$, which is a model
extent and not a stature. The model is posed with the first row of the
coordinate matrix before the extent is read, so $h$ does not depend on the
pose in which the run ends or on the model's up axis. The source of $h$ is
recorded as \texttt{body\_height\_source} (\texttt{body\_height\_m} or
\texttt{model\_extent\_first\_frame}) in the residual budget and, with $h$, in
the study receipt (\texttt{id\_body\_height\_m}, \texttt{id\_body\_height\_source};
receipt schema \texttt{spine-study-receipt/2}). Without external loads the residual is the
total weight, so the budget fails by construction.

\paragraph{External Loads.}
An ExternalLoads XML is added to the model with \texttt{ModOpAddExternalLoads}
and read at the state time of each frame. The provenance key
\texttt{external\_loads} states exactly what was applied, or reads
\texttt{none (excluded): gravity only}.

\paragraph{Derivatives.}
Coordinates are optionally low-passed (zero-phase Butterworth,
\texttt{smooth\_kinematic\_trajectory}) and differentiated with a not-a-knot
cubic spline evaluated at the sample times. Pelvis Euler rates map to
ground-frame angular velocity and acceleration by
$\omega = J\dot q$ and
$\alpha = J\ddot q + \dot q_2\,\dot c_2 + \dot q_3\,\dot c_3$, with
$\dot c_2=\dot q_1(c_1\times c_2)$ and
$\dot c_3=\dot q_1(c_1\times c_3)+\dot q_2(c_2\times c_3)$. Spline end
effects make the first and last samples less reliable.

\paragraph{Assumptions and Limitations.}
Pelvis parent and child offset frames are identity. The level axes are
ISB-aligned to the inventory tolerance. The fixture parameters are invented;
no value is transcribed from a licensed model. Residual verdicts are software
checks, not scientific qualification.

\paragraph{Verification and Reproduction.}
The tests compare the static level moments and the pelvis wrench with an
independent numpy rigid-body composition to $10^{-6}$~N$\cdot$m, and a dynamic
motion with $M(q)\ddot q + C + G$ from the matter subsystem.
\begin{verbatim}
python -m pytest tests/opensim/spine/test_inverse_dynamics.py \
    -q -n 0 -p no:cacheprovider -o addopts=""
\end{verbatim}
